\newpage
\chapter{OASIS3-MCT C Bindings}
\label{sec_cbindings}

OASIS3-MCT is distributed with C bindings and can be called from models written
in C and in C++.  These bindings leverage the Fortran ISO\_C\_BINDING standard.  The model
interface names match, more or less, the Fortran implementation and are documented
below.  The cbindings can be compiled into static or shared libraries by the
OASIS3-MCT TopMakefileOasis3 as documented in \ref{subsec_compile}.

To use the OASIS3-MCT C bindings, the statement

\begin{verbatim}
#include "oasis_c.h"
\end{verbatim}

needs to be added to the C model source code.  All available
parameters macros and
interfaces are defined there.


\section{Model interfaces for C}
\label{sec_cinterfaces}

The C interfaces are specified below.  These interfaces largely match
up with equivalent interfaces in Fortran.  An interface named
{\tt oasis\_interface} in Fortran can be expected to be named
{\tt oasis\_c\_interface} in C.

All of the C interfaces return an integer error code which can be
tested against the {\tt OASIS\_Ok} (or the equivalent {\tt
  OASIS\_Success}) constant.\\
The {\tt OASIS\_CHECK\_ERR} macro aborts
OASIS with a meaningful message in the debug files in case of failure.

For most of the functions,
the return code is consistent with the {\tt
  kinfo} argument in the Fortran interfaces.  For these functions, the
following constructs are equivalent in the two languages:

\begin{verbatim}
call oasis_get_localcomm(localcomm, kinfo)
if (kinfo .ne. OASIS_Ok)  &
   &  call OASIS_Abort(comp_id, "oasis_get_localcomm", &
   &  "Runtime error", __FILE__, __LINE__)
\end{verbatim}

and

\begin{verbatim}
OASIS_CHECK_ERR(oasis_get_localcomm(localcomm))
\end{verbatim}

For convenience a similar macro {\tt OASIS\_CHECK\_MPI\_ERR} has been
defined for testing the return code of any MPI function against {\tt
  MPI\_SUCCESS} and cleanly aborting OASIS3-MCT in case of failure.\\

For the communication functions, the return error code only indicates
success or failure, while the detailed status is returned by the {\tt
  kinfo} argument.

\vspace{0.2cm}
More information about all these interfaces can be found in the
Fortran interface section (see \ref{API}).  Notes and deviations from
the Fortran standard are noted below.\\
Use of these C bindings are illustrated in practical examples in
directories {\tt /C} in the different subdirectories in {\tt pyoasis/examples/}

\vspace{0.2cm}
\begin{itemize}

\item {\tt int oasis\_c\_init\_comp(int* compid, const char* comp\_name, const bool coupled)}
  \begin{itemize}
  \item This {\tt oasis\_c\_init\_comp} interface includes the {\tt coupled} flag but
    NO communicator argument ({\tt commworld}); to indicate a
    communicator argument, see {\tt oasis\_c\_init\_comp\_with\_comm} below.
    \item The {\tt coupled} argument can receive the predefined mnemonic boolean constants {\tt
      OASIS\_Coupled} and {\tt OASIS\_Not\_Coupled}
  \end{itemize}

\vspace{0.2cm}
\item {\tt int oasis\_c\_init\_comp\_with\_comm(int* compid, const char* comp\_name, const bool coupled, const MPI\_Comm commworld)}
  \begin{itemize}
  \item This alternative {\tt init\_comp} interface has the {\tt coupled} flag and a C {\tt MPI\_Comm} communicator
  \end{itemize}

\vspace{0.2cm}
\item {\tt int oasis\_c\_get\_localcomm(MPI\_Comm* local\_comm)}
  \begin{itemize}
  \item Returns the local communicator by reference in a C {\tt MPI\_Comm*} type argument
  \end{itemize}

\vspace{0.2cm}
\item {\tt int oasis\_c\_create\_couplcomm(const int icpl, const MPI\_Comm local\_comm, MPI\_Comm* coupl\_comm)}
  \begin{itemize}
  \item {\tt loca\_comm} is a C {\tt MPI\_Comm} type input argument (by value)
  \item {\tt coupl\_comm} is a C {\tt MPI\_Comm*} type output argument (by reference)
  \end{itemize}

\vspace{0.2cm}
\item {\tt int oasis\_c\_set\_couplcomm(const MPI\_Comm coupl\_comm)}
  \begin{itemize}
  \item Takes a C {\tt MPI\_Comm} type as an input argument (by value)
  \end{itemize}

\vspace{0.2cm}
\item {\tt int oasis\_c\_def\_partition(int* il\_part\_id, const int ig\_paral\_size, const int* ig\_paral, const int ig\_size, const char* name)}
  \begin{itemize}
  \item {\tt ig\_paral\_size} is the size of the array {\tt ig\_paral}: the
    {\tt oasis\_c.h} header provides a set of constants and macros for
    the size of every partition strategy, namely {\tt
      OASIS\_Serial\_Params}, {\tt OASIS\_Apple\_Params}, {\tt
      OASIS\_Box\_Params}, {\tt
      OASIS\_Orange\_Params(n\_segments)}, {\tt
      OASIS\_Points\_Params(n\_points)}
  \item the {\tt oasis\_c.h} header also provides a set of predefined
    constants for the storages positions in {\tt
    ig\_paral}, namely {\tt OASIS\_Strategy}, {\tt OASIS\_Segments}, {\tt
    OASIS\_Npoints}, {\tt OASIS\_Offset}, {\tt OASIS\_Length}, {\tt
    OASIS\_SizeX}, {\tt OASIS\_SizeY}, {\tt OASIS\_LdX}
  \item the partition strategy, to be stored in the {\tt
    OASIS\_Strategy} position of {\tt ig\_paral}, can take one of the
    following predefined constants values: {\tt OASIS\_Serial}, {\tt
      OASIS\_Apple}, {\tt OASIS\_Box}, {\tt OASIS\_Orange}, {\tt
      OASIS\_Points}
  \item for the cases in which the {\tt ig\_size} and {\tt name}
    arguments are
    not relevant for the partition definition, the placeholders {\tt
      OASIS\_No\_Gsize} and {\tt OASIS\_No\_Name} can be used instead
  \end{itemize}
  Here an example for the {\tt OASIS\_Part\_Apple} strategy
\begin{verbatim}
  int part_params[OASIS_Apple_Params];
  part_params[OASIS_Strategy] = OASIS_Apple;
  part_params[OASIS_Offset] = offset;
  part_params[OASIS_Length] = local_size;
  int part_id;
  OASIS_CHECK_ERR(oasis_c_def_partition(&part_id, OASIS_Apple_Params,
                                        part_params, OASIS_No_Gsize,
                                        OASIS_No_Name));
\end{verbatim}

\vspace{0.2cm}
\item {\tt int oasis\_c\_start\_grids\_writing()}

\vspace{0.2cm}
\item {\tt int oasis\_c\_write\_grid(const char* cgrid, const int nx\_global, const int ny\_global, const int nx\_loc, const int ny\_loc, const double* lon, const double* lat, const int il\_partid)}
  \begin{itemize}
    \item {\tt nx\_global} and {\tt ny\_global} are the first and second dimensions
      of the global grid
    \item {\tt nx\_loc} and {\tt ny\_loc} are the two dimensions
      of the local arrays {\tt lon} and {\tt lat}
    \item {\tt lon} and {\tt lat} are stored with a fortran compatible
      (column major) ordering. If they are declared as double
      pointers, they have to be {\tt lon[ny\_loc][nx\_loc]}
    \item {\tt il\_partid} is relevant only for parallel writing. In case
      of single proc invocations, the constant {\tt OASIS\_No\_Part}
      can be passed as a placeholder
  \end{itemize}

\vspace{0.2cm}
\item {\tt int oasis\_c\_write\_corner(const char* cgrid, const int nx\_global, const int ny\_global, const int nc, const int nx\_loc, const int ny\_loc, const double* clon, const double* clat, const int il\_partid)}
  \begin{itemize}
    \item {\tt nx\_global} and {\tt ny\_global} are the first and second dimensions
      of the global grid
    \item {\tt nc} is the maximum number of corners per cell
    \item {\tt nx\_loc} and {\tt ny\_loc} are the two dimensions
      of the local arrays {\tt clon} and {\tt clat}
    \item {\tt clon} and {\tt clat} are stored with a fortran compatible
      (column major) ordering. If they are declared as triple
      pointers, they have to be {\tt clon[nc][ny\_loc][nx\_loc]}
    \item {\tt il\_partid} is relevant only for parallel writing. In case
      of single proc invocations, the constant {\tt OASIS\_No\_Part}
      can be passed as a placeholder
  \end{itemize}

\vspace{0.2cm}
\item {\tt int oasis\_c\_write\_mask(const char* cgrid, const int nx\_global, const int ny\_global, const int nx\_loc, const int ny\_loc, const int* mask, const int il\_partid, const char* companion)}
  \begin{itemize}
    \item {\tt nx\_global} and {\tt ny\_global} are the first and second dimensions
      of the global grid
    \item {\tt nx\_loc} and {\tt ny\_loc} are the two dimensions
      of the local array {\tt mask}
    \item {\tt mask} is stored with a fortran compatible
      (column major) ordering. If it is  declared as a double
      pointer, it has to be {\tt mask[ny\_loc][nx\_loc]}
    \item {\tt il\_partid} is relevant only for parallel writing. In case
      of single proc invocations, the constant {\tt OASIS\_No\_Part}
      can be passed as a placeholder
    \item if no {\tt companion} grid attribute is needed, the constant
      {\tt OASIS\_No\_Companion} can be passed as a placeholder
  \end{itemize}

\vspace{0.2cm}
\item {\tt int oasis\_c\_write\_frac(const char* cgrid, const int nx\_global, const int ny\_global, const int nx\_loc, const int ny\_loc, const double* frac, const int il\_partid, const char* companion)}
  \begin{itemize}
    \item {\tt nx\_global} and {\tt ny\_global} are the first and second dimensions
      of the global grid
    \item {\tt nx\_loc} and {\tt ny\_loc} are the two dimensions
      of the local array {\tt frac}
    \item {\tt frac} is stored with a fortran compatible
      (column major) ordering. If it is  declared as a double
      pointer, it has to be {\tt frac[ny\_loc][nx\_loc]}
    \item {\tt il\_partid} is relevant only for parallel writing. In case
      of single proc invocations, the constant {\tt OASIS\_No\_Part}
      can be passed as a placeholder
    \item if no {\tt companion} grid attribute is needed, the constant
      {\tt OASIS\_No\_Companion} can be passed as a placeholder
  \end{itemize}

\vspace{0.2cm}
\item {\tt int oasis\_c\_write\_area(const char* cgrid, const int nx\_global, const int ny\_global, const int nx\_loc, const int ny\_loc, const double* area, const int il\_partid)}
  \begin{itemize}
    \item {\tt nx\_global} and {\tt ny\_global} are the first and second dimensions
      of the global grid
    \item {\tt nx\_loc} and {\tt ny\_loc} are the two dimensions
      of the local array {\tt area}
    \item {\tt area} is stored with a fortran compatible
      (column major) ordering. If it is  declared as a double
      pointer, it has to be {\tt area[ny\_loc][nx\_loc]}
    \item {\tt il\_partid} is relevant only for parallel writing. In case
      of single proc invocations, the constant {\tt OASIS\_No\_Part}
      can be passed as a placeholder
  \end{itemize}

\vspace{0.2cm}
\item {\tt int oasis\_c\_write\_angle(const char* cgrid, const int nx\_global, const int ny\_global, const int nx\_loc, const int ny\_loc, const double* angle, const int il\_partid)}
  \begin{itemize}
    \item {\tt nx\_global} and {\tt ny\_global} are the first and second dimensions
      of the global grid
    \item {\tt nx\_loc} and {\tt ny\_loc} are the two dimensions
      of the local array {\tt angle}
    \item {\tt angle} is stored with a fortran compatible
      (column major) ordering. If it is  declared as a double
      pointer, it has to be {\tt angle[ny\_loc][nx\_loc]}
    \item {\tt il\_partid} is relevant only for parallel writing. In case
      of single proc invocations, the constant {\tt OASIS\_No\_Part}
      can be passed as a placeholder
  \end{itemize}

\vspace{0.2cm}
\item {\tt int oasis\_c\_terminate\_grids\_writing()}

\vspace{0.2cm}
\item {\tt int oasis\_c\_def\_var(int* var\_id, const char* name, const int il\_part\_id, const bundle\_size, const int kinout, const int var\_type)}
  \begin{itemize}
  \item {\tt bundle\_size} is a scalar integer, corresponding to the
    second entry of the fortran {\tt var\_nodims} array
  \item {\tt kinout} can receive one of the two predefined constants
    {\tt OASIS\_In} or {\tt OASIS\_Out} (also in the form
    {\tt OASIS\_IN} or {\tt OASIS\_OUT})
  \item {\tt var\_type} can receive one of the two predefined constants
    {\tt OASIS\_Real} or {\tt OASIS\_Double} (also in the form
    {\tt OASIS\_REAL} or {\tt OASIS\_DOUBLE})
  \end{itemize}

\vspace{0.2cm}
\item {\tt int oasis\_c\_enddef()}

\vspace{0.2cm}
\item {\tt int oasis\_c\_put(const int var\_id, const int date, const
  int x\_size, const int y\_size, const int bundle\_size, const int
  fkind, const int storage, const void* fld1, const bool
  write\_restart, int* kinfo)}
  \begin{itemize}
  \item This interface does not support higher order mapping through optional fields at this time.
  \item {\tt x\_size}, {\tt y\_size}, are the dimensions of the local
    portion of the domain in the order as they appear in the {\tt
      grids.nc} file and which corresponds to the storage order in the
    corresponding internal OASIS Fortran work arrays. \\
    Notice that for
    unstructured grids {\tt y\_size=1}
  \item {\tt bundle\_size} should be coherent with the argument of
    same name in the corresponding {\tt oasis\_c\_def\_var}
  \item {\tt fkind} can take one of the two predefined constants
    {\tt OASIS\_Real} or {\tt OASIS\_Double} (also in the form
    {\tt OASIS\_REAL} or {\tt OASIS\_DOUBLE}) should be coherent with
    the argument {\tt ktype} in {\tt oasis\_c\_def\_var}
  \item {\tt storage} can take one of the two predefined constants
    {\tt OASIS\_COL\_MAJOR} or \\
    {\tt OASIS\_ROW\_MAJOR}. \\
    In the first
    case, the array containing the field has to be declared as \\ {\tt
      field[bundle\_size][y\_size][x\_size]} (or any equivalent
    storage of total size \\
    {\tt bundle\_size*y\_size*x\_size} stored in
    column major order. \\
    In the second case, the field has to be
    declared as \\
    {\tt field[x\_size][y\_size][bundle\_size]} (or any
    equivalent row major storage). \\
    Notice that this choice implies an
    extra memory copy performed internally by OASIS before acting on
    the field and is therefore to be avoided whenever possible.
  \item {\tt write\_restart} can take one of the two predefined constants
    {\tt OASIS\_Write\_Restart} or {\tt OASIS\_No\_Restart}
  \item {\tt kinfo} contains a return status to be compared against
    the values of the {\tt return\_codes} enumeration in {\tt oasis\_c.h}
  \end{itemize}

\vspace{0.2cm}
\item {\tt int oasis\_c\_get(const int var\_id, const int date, const
  int x\_size, const int y\_size, const int bundle\_size, const int
  fkind, const int storage, void* fld1, int* kinfo)}
  \begin{itemize}
  \item arguments are the same than for {\tt oasis\_c\_put}.
  \end{itemize}

\vspace{0.2cm}
\item {\tt int oasis\_c\_terminate()}

\vspace{0.2cm}
\item {\tt int oasis\_c\_abort(const int compid, const char* routine\_name, const char* abort\_message, const char* file, const int line)}

\vspace{0.2cm}
\item {\tt int oasis\_c\_get\_debug(int* debug)}

\vspace{0.2cm}
\item {\tt int oasis\_c\_set\_debug(const int debug)}

\vspace{0.2cm}
\item {\tt int oasis\_c\_get\_intercomm(MPI\_Comm* new\_comm, char* cdnam)}
  \begin{itemize}
  \item Returns a C MPI\_Comm type by reference
  \end{itemize}

\vspace{0.2cm}
\item {\tt int oasis\_c\_get\_intracomm(MPI\_Comm* new\_comm, char* cdnam)}
  \begin{itemize}
  \item Returns a C MPI\_Comm type as by reference
  \end{itemize}

\vspace{0.2cm}
\item {\tt int oasis\_c\_get\_multi\_intracomm(MPI\_Comm* new\_comm, const int cdnam\_size, char** cdnam, int* root\_ranks)}
  \begin{itemize}
  \item Returns a C MPI\_Comm type as by reference
  \item cdnam\_size is the size of the array cdnam
  \end{itemize}

\vspace{0.2cm}
\item {\tt int oasis\_c\_put\_inquire(int var\_id, int date, int* kinfo)}
  \begin{itemize}
  \item {\tt kinfo} contains a return status to be compared against
    the values of the {\tt return\_codes} enumeration in {\tt oasis\_c.h}
  \end{itemize}

\vspace{0.2cm}
\item {\tt int oasis\_c\_get\_ncpl(const int varid, int* ncpl)}

\vspace{0.2cm}
\item {\tt int oasis\_c\_get\_freqs(const int varid, const int mop, const int ncpl, int*  cpl\_freqs)}

\end{itemize}
